% ============================================================================
% Fichier  : partie4/P04_C16_reseaux_equipements.tex
% Titre    : Sécurité des Réseaux (2) : Équipements et Protocoles Sûrs
% Auteur   : MINKA MI NGUIDJOÏ Thierry Emmanuel
% Version  : 1.0 -- novembre 2026
% ============================================================================
% Base : brouillon « Livre V » (périphériques, communications sécurisées).
% Sources : NIST SP 800-41 Rev. 1 (pare-feu) ; NIST SP 800-94 (IDPS) ;
% RFC 7296 (IKEv2) ; Vanhoef et Piessens (2017, KRACK) ; IEEE 802.1X.
% ============================================================================

\chapter{Sécurité des Réseaux (2) : Équipements et Protocoles Sûrs}
\label{chap:reseaux2}
\chaptermark{Équipements et protocoles sûrs}

\epigraph{%
    « Un pare-feu applique une politique~;\\
    il ne la remplace pas. »%
}{Principe du cours}

\niveauChapitre{Intermédiaire}{Chapitres~\ref{chap:architecture}
et~\ref{chap:reseaux1}.}

\begin{boxobjectifs}
\begin{itemize}
    \item Distinguer les types de pare-feu et lire une table de règles.
    \item Repérer les règles trop permissives, ombrées ou inutiles.
    \item Expliquer le fonctionnement et les limites des IDS et IPS.
    \item Choisir un réseau privé virtuel (IPsec, TLS, WireGuard).
    \item Sécuriser un réseau Wi-Fi d'entreprise.
\end{itemize}
\end{boxobjectifs}

\bigskip

% ============================================================================
\section*{Question directrice : les pare-feu de LiZbiZ servaient-ils à quelque chose~?}
% ============================================================================

LiZbiZ avait un pare-feu dans chaque centre. Lors de l'enquête, la revue
des règles de Garoua montre que la première règle autorisait « tout
trafic » depuis le réseau du prestataire, et qu'une règle plus précise,
écrite plus bas, n'était jamais appliquée. Un pare-feu mal réglé donne
une impression de sécurité sans la sécurité.

\medskip
\textbf{Plan du chapitre.}
\begin{enumerate}
    \item Les pare-feu (section~\ref{sec:c16-parefeu}).
    \item Auditer une table de règles (section~\ref{sec:c16-regles}).
    \item Détecter et prévenir les intrusions (section~\ref{sec:c16-ids}).
    \item Les réseaux privés virtuels (section~\ref{sec:c16-vpn}).
    \item Le Wi-Fi (section~\ref{sec:c16-wifi}).
\end{enumerate}

\separateur

% ============================================================================
\section{Les Pare-feu}
\label{sec:c16-parefeu}
% ============================================================================

\begin{tableaucours}{L{3.2cm}L{5.0cm}L{4.0cm}}{%
    \tableauHeader{Type} &
    \tableauHeader{Ce qu'il examine} &
    \tableauHeader{Limite}
}
Filtrage sans état & Chaque paquet isolément~: adresses, ports, protocole
& Ne sait pas si un paquet appartient à une connexion légitime \\
Filtrage avec état & Les connexions~: n'accepte une réponse que si la
demande est sortie & Ne voit pas le contenu applicatif \\
Pare-feu applicatif, mandataire & Le contenu d'un protocole (HTTP,
DNS) & Plus lent~; limité aux protocoles qu'il connaît \\
Pare-feu de nouvelle génération & État, applications, utilisateurs,
réputation, parfois inspection du chiffré & Complexité, coût \\
\end{tableaucours}
\vspace{-0.8em}
\begin{center}{\small\itshape Tableau~16.1. Les types de pare-feu~\cite{nist80041}}\end{center}

Un pare-feu lit ses règles \emph{dans l'ordre} et applique la première qui
correspond. La dernière règle est un refus général~: tout ce qui n'est
pas explicitement autorisé est interdit (valeur par défaut sûre,
Chapitre~\ref{chap:proprietes}).

% ============================================================================
\section{Auditer une Table de Règles}
\label{sec:c16-regles}
% ============================================================================

\begin{tableaucours}{cL{2.6cm}L{2.6cm}L{1.8cm}L{1.6cm}L{2.6cm}}{%
    \tableauHeader{N°} &
    \tableauHeader{Source} &
    \tableauHeader{Destination} &
    \tableauHeader{Port} &
    \tableauHeader{Action} &
    \tableauHeader{Commentaire}
}
1 & Réseau prestataire & Tout & Tout & Autoriser & « temporaire, 2026 » \\
2 & Réseau prestataire & Serveurs Garoua & 22, 3389 & Autoriser & Maintenance \\
3 & Postes & Serveur APP-03 & 443 & Autoriser & APP-03 retiré en 2025 \\
4 & Tout & Portail & 443 & Autoriser & Portail de suivi \\
5 & Postes & Internet & Tout & Autoriser & Navigation \\
6 & Tout & Tout & Tout & Refuser & Règle finale \\
\end{tableaucours}
\vspace{-0.8em}
\begin{center}{\small\itshape Tableau~16.2. Règles du pare-feu de Garoua (extrait,
mars)}\end{center}

\begin{boxscenario}{Opération GAROUA. La revue des règles}
\textbf{Constatations.} La règle~1 autorise tout depuis le réseau du
prestataire~: elle \emph{ombre} la règle~2, qui ne s'applique jamais. La
règle~3 vise un serveur retiré~: elle est inutile. La règle~5 autorise
tous les ports vers Internet depuis les postes~: c'est par elle que sont
sortis les 40~Go.

\medskip
\textbf{Analyse} (confiance élevée). Supprimer la règle~1, limiter la
règle~2 aux serveurs et aux heures de maintenance, supprimer la règle~3,
restreindre la règle~5 aux ports web via un mandataire filtrant. Chaque
règle doit avoir un propriétaire et une date de revue.
\end{boxscenario}

\begin{boxprocedure}
\textbf{Revoir une table de règles}
\begin{enumerate}
    \item Repérer les règles « tout » (source, destination ou port).
    \item Repérer les règles ombrées par une règle plus large située
          au-dessus.
    \item Repérer les règles qui visent des systèmes retirés, ou jamais
          utilisées (compteurs à zéro).
    \item Vérifier que chaque règle a un motif, un propriétaire et une date.
    \item Vérifier la règle finale de refus et la journalisation des refus.
\end{enumerate}
\end{boxprocedure}

% ============================================================================
\section{Détecter et Prévenir les Intrusions}
\label{sec:c16-ids}
% ============================================================================

\begin{boxdef}
\termecle{IDS} (système de détection d'intrusion)~: analyse le trafic ou
l'activité d'une machine et émet des alertes. \termecle{IPS} (système de
prévention)~: placé en coupure, il bloque ce qu'il détecte. La détection
repose sur des \emph{signatures} (motifs d'attaques connues) ou sur des
\emph{anomalies} (écarts au comportement habituel)~\cite{nist80094}.
\end{boxdef}

\begin{tableaucours}{L{3.4cm}L{4.4cm}L{4.4cm}}{%
    \tableauHeader{Méthode} &
    \tableauHeader{Forces} &
    \tableauHeader{Faiblesses}
}
Signatures & Précise sur les attaques connues~; peu de faux positifs &
Aveugle aux attaques nouvelles et au trafic chiffré \\
Anomalies & Peut voir l'inconnu & Faux positifs nombreux~; apprentissage
long \\
\end{tableaucours}
\vspace{-0.8em}
\begin{center}{\small\itshape Tableau~16.3. Signatures et anomalies}\end{center}

\begin{boxalert}
\textbf{Violation typique.} Installer un IDS et ne jamais lire ses alertes.

\smallskip
\textbf{Règle~:} un détecteur sans personne pour traiter ses alertes ne
détecte rien. On règle les signatures pour réduire les faux positifs, et
on organise le traitement des alertes (Chapitre~\ref{chap:operations}).
\end{boxalert}

% ============================================================================
\section{Les Réseaux Privés Virtuels}
\label{sec:c16-vpn}
% ============================================================================

Un \termecle{réseau privé virtuel} (VPN) crée un tunnel chiffré et
authentifié entre deux réseaux, ou entre un poste et un réseau.

\begin{tableaucours}{L{2.6cm}L{5.0cm}L{4.6cm}}{%
    \tableauHeader{Technologie} &
    \tableauHeader{Principe} &
    \tableauHeader{Usage}
}
IPsec (IKEv2) & Chiffrement au niveau IP, négociation des clés par
IKEv2~\cite{rfc7296} & Liaisons entre sites (Douala et Garoua) \\
VPN TLS & Tunnel au-dessus de TLS & Accès distant des nomades \\
WireGuard & Protocole récent, code réduit, cryptographie moderne & Accès
distant et liaisons entre sites \\
\end{tableaucours}
\vspace{-0.8em}
\begin{center}{\small\itshape Tableau~16.4. Les réseaux privés virtuels}\end{center}

Un VPN protège le transport~; il ne vérifie ni la santé du poste ni la
légitimité de l'utilisateur au-delà de son authentification. C'est
pourquoi l'accès distant exige un second facteur et ne doit ouvrir que les
ressources nécessaires (Chapitre~\ref{chap:architecture}).

% ============================================================================
\section{Le Wi-Fi}
\label{sec:c16-wifi}
% ============================================================================

\begin{tableaucours}{L{2.4cm}L{9.8cm}}{%
    \tableauHeader{Protection} &
    \tableauHeader{Statut}
}
WEP & Cassé~; à proscrire \\
WPA2-Personnel & Clé partagée~; vulnérable aux dictionnaires si la phrase
est faible~; la vulnérabilité KRACK (2017) a exigé des
correctifs~\cite{vanhoef2017} \\
WPA3-Personnel & Authentification SAE, résistante aux attaques par
dictionnaire hors ligne \\
WPA2/WPA3-Entreprise & Authentification individuelle par 802.1X et un
serveur RADIUS~; recommandée en entreprise \\
\end{tableaucours}
\vspace{-0.8em}
\begin{center}{\small\itshape Tableau~16.5. Les protections du Wi-Fi}\end{center}

\begin{boxmetier}{Administrateur systèmes et réseaux}
\textbf{Faire vivre les règles.} L'administrateur revoit les règles des
pare-feu chaque trimestre, supprime ce qui n'a plus de motif, et
journalise les refus. Il sépare le Wi-Fi des visiteurs du réseau interne.
\end{boxmetier}

\begin{boxmetier}{Analyste SOC}
\textbf{Régler la détection.} L'analyste ajuste les signatures de l'IDS au
contexte (désactive celles qui visent des systèmes absents, crée des
règles propres à l'entreprise) et suit le taux de faux positifs.
\end{boxmetier}

\separateur

\begin{boxresume}
\textbf{Ce qu'il faut retenir de ce chapitre~:}
\begin{itemize}
    \item \textbf{Les pare-feu} filtrent sans état, avec état, ou au niveau
          applicatif~; les règles se lisent dans l'ordre, avec un refus
          final.
    \item \textbf{Une table de règles s'audite}~: règles « tout », ombrées,
          inutiles, sans propriétaire.
    \item \textbf{IDS et IPS} détectent par signatures ou anomalies~; sans
          traitement des alertes, ils ne servent à rien.
    \item \textbf{Les VPN} (IPsec, TLS, WireGuard) protègent le transport,
          pas la légitimité de l'accès.
    \item \textbf{Le Wi-Fi d'entreprise} utilise WPA3 ou 802.1X~; WEP est
          à proscrire.
\end{itemize}
\end{boxresume}

\begin{boxinfo}
\textbf{Points de contrôle de l'élève.} Je sais~: distinguer filtrage sans
état, avec état et applicatif~; repérer une règle ombrée~; conduire la
revue d'une table de règles~; comparer signatures et anomalies~; choisir
un VPN et une protection Wi-Fi.
\end{boxinfo}

\bigskip

\begin{boxexo}[title={\ding{45}\quad Exercice~16.1. Avec ou sans état \niveaubase}]
Pourquoi un filtrage avec état permet-il de n'écrire qu'une règle pour la
navigation web des postes, là où un filtrage sans état en demande deux~?
\end{boxexo}

\begin{boxexo}[title={\ding{45}\quad Exercice~16.2. L'ordre \niveaubase}]
Dans le Tableau~16.2, quelle règle s'applique à une connexion du réseau
du prestataire vers le contrôleur de domaine, port 445~? Et après
suppression de la règle~1~?
\end{boxexo}

\begin{boxexo}[title={\ding{45}\quad Exercice~16.3. Wi-Fi \niveaubase}]
Une agence utilise encore WEP pour son Wi-Fi interne. Que recommandez-vous,
et pourquoi~?
\end{boxexo}

\begin{boxexo}[title={\ding{45}\quad Exercice~16.4. Réécrire les règles \niveaumoyen}]
Réécrivez la table du Tableau~16.2 après la revue, avec motifs,
propriétaires et dates.
\end{boxexo}

\begin{boxexo}[title={\ding{45}\quad Exercice~16.5. Faux positifs \niveaumoyen}]
L'IDS de la BANQUE SANAGA produit 4\,000~alertes par jour, dont 98~\% de
faux positifs. Combien de vraies alertes par jour~? Que faire pour que
l'équipe puisse les traiter~?
\end{boxexo}

\begin{boxexo}[title={\ding{45}\quad Exercice~16.6. IPS ou IDS \niveaumoyen}]
Faut-il placer un IPS en coupure devant le système central de la BANQUE
SANAGA~? Pesez le risque d'un blocage à tort.
\end{boxexo}

\begin{boxexo}[title={\ding{45}\quad Exercice~16.7. VPN \niveaumoyen}]
Choisissez le VPN pour~: relier Douala et Garoua~; donner l'accès aux
commerciaux nomades~; donner l'accès au prestataire. Justifiez.
\end{boxexo}

\begin{boxexo}[title={\ding{45}\quad Exercice~16.8. Le Wi-Fi des visiteurs \niveaumoyen}]
Concevez le Wi-Fi d'une agence~: réseau du personnel et réseau des
visiteurs, authentification, séparation, filtrage.
\end{boxexo}

\begin{boxexo}[title={\ding{45}\quad Exercice~16.9. Le filtrage sortant \niveauexpert}]
Proposez une politique de filtrage sortant pour LiZbiZ qui aurait arrêté
ou signalé l'exfiltration des 40~Go, sans bloquer l'activité.
\end{boxexo}

\begin{boxexo}[title={\ding{45}\quad Exercice~16.10. Un programme de revue \niveauexpert}]
Rédigez la procédure trimestrielle de revue des pare-feu de LiZbiZ et
l'indicateur qui en rend compte au comité de sécurité.
\end{boxexo}

\bigskip

\begin{boxinfo}
\textbf{Pour aller plus loin.}
\begin{itemize}
    \item NIST SP~800-41 (pare-feu) et SP~800-94 (détection
          d'intrusion)~\cite{nist80041,nist80094}.
    \item RFC~7296 (IKEv2)~\cite{rfc7296}~; M.~Vanhoef et F.~Piessens sur
          KRACK~\cite{vanhoef2017}.
\end{itemize}
\end{boxinfo}

\begin{boxinfo}
\textbf{Transition.} Le réseau filtre les flux~; il faut encore savoir qui
les émet. Le Chapitre~\ref{chap:iam} traite des identités et des accès.
\end{boxinfo}
